\section{What binary structure preserves}
For binary features, write $q_j=\Prob(X_j=1)$ and $p_{11}=\Prob(X_1=X_2=1)$,
or use their empirical counterparts for a fixed design. Then
\[
 a=q_1(1-q_1),\quad b=q_2(1-q_2),\quad c=p_{11}-q_1q_2,
 \qquad \max(0,q_1+q_2-1)\le p_{11}\le\min(q_1,q_2).
\]
The conditional mean identities
\begin{equation}\label{eq:conditional-differences}
 \frac ca=\E(X_2\mid X_1=1)-\E(X_2\mid X_1=0),\qquad
 \frac cb=\E(X_1\mid X_2=1)-\E(X_1\mid X_2=0)
\end{equation}
give $|c|\le\min(a,b)$. Equality in either bound forces identical or
complementary columns, so both inequalities are strict when $C$ is positive
definite.

\begin{theorem}[Exact single-signal shrinkage]\label{thm:single}
In the noiseless two-binary-feature lasso, let $u=(a,b)$, $r=0$, and
$\beta^*=(B,0)$. For every $t\ge0$, the unique minimizer is
$\hat\beta=(\soft_t(B),0)$.
\end{theorem}
\begin{proof}
The first coordinate satisfies its scalar KKT condition. The inactive residual
has absolute value $|c(B-\soft_t(B))|\le |c|t\le bt$, so the second KKT
condition also holds. Strict convexity gives uniqueness.
\end{proof}
The result also holds with coordinates exchanged. Noise can still produce
false selections, so this is a statement about noiseless shrinkage.

For two active coefficients with signs $h$, \eqref{eq:both-active} gives
\begin{equation}\label{eq:binary-shrinkage}
 \hat\beta_1=\beta_1^*-\frac{t\{h_1-(c/a)h_2\}}{1-\rho^2},\qquad
 \hat\beta_2=\beta_2^*-\frac{t\{h_2-(c/b)h_1\}}{1-\rho^2},\qquad
 \rho=\frac{c}{\sqrt{ab}}.
\end{equation}
The resulting coefficients must have the stipulated signs. For a nested
design with $q_1=1/2$, $q_2=p_{11}=1/10$, we have
$(a,b,c)=(.25,.09,.05)$. At $t=.1$, the exact solutions are
\begin{center}
\begin{tabular}{lll}
\toprule
$\beta^*$ & Orthogonal, same marginals & Nested design \\
\midrule
$(1,1)$ & $(.9,.9)$ & $(.91,.95)$ \\
$(1,-1)$ & $(.9,-.9)$ & $(.865,-.825)$ \\
$(1,0)$ & $(.9,0)$ & $(.9,0)$ \\
\bottomrule
\end{tabular}
\end{center}
Thus correlation can increase or decrease shrinkage depending on the signs.
Both designs can be realized with integer cell counts, using sample sizes
that are multiples of twenty for the orthogonal comparison.

\begin{proposition}[No universal diagonal correction]\label{prop:no-diagonal}
If $c\ne0$ and $t>0$, no positive design-dependent weights $u$ make the
noiseless lasso equal $(\soft_t(\beta_1^*),\soft_t(\beta_2^*))$ for every
$\beta^*\in\R^2$.
\end{proposition}
\begin{proof}
For $(B,0)$ with $B>t$, the desired active KKT equation forces $u_1=a$.
Exchanging coordinates forces $u_2=b$. For $(B,B)$ with $B>t$, the desired
solution would then require $C(1,1)^\top=(a,b)^\top$, which forces $c=0$.
\end{proof}
If active signs were known, the weights $u_1=a+ch_1h_2$ and $u_2=b+ch_1h_2$
would instead produce $\hat\beta=\beta^*-th$ whenever $\beta^*$ has signs $h$
and $\min_j|\beta_j^*|>t$. Binary feasibility makes these weights positive.
Their dependence on unknown signs prevents a universal preprocessing rule.

\begin{proposition}[An approximation condition]\label{prop:bound}
Let $D=\diag(a,b)$ and $\eta=\max(|c|/a,|c|/b)<1$. On a valid noiseless
both-active sign region with variance weights,
\[
 \|\hat\beta-(\beta^*-th)\|_\infty\le t\frac{\eta}{1-\eta}.
\]
\end{proposition}
\begin{proof}
Write $C=D(I+F)$, where $\|F\|_\infty=\eta$. The Neumann series gives
$\|(I+F)^{-1}-I\|_\infty\le\eta/(1-\eta)$, and
$\hat\beta=\beta^*-tC^{-1}Dh$.
\end{proof}
Small conditional mean differences in \eqref{eq:conditional-differences}
therefore control the approximation. Small Pearson correlation alone does
not provide this control when variances differ greatly.
